\SourceLeaf{151}Zusammenhang und erkenne darin nichts weiter als eine bloße Petitio principii. Sie werden mich richtig verstehen. Ich lege wenig Wert auf eine streng logische Formulierung, die in Schriften für Männer und Kenner oft nur Pedanterie wäre. Der streng logische Zusammenhang muss sich aber überall nachweisen lassen, wenn es verlangt wird. Das ist eine unerlässliche Bedingung eines guten Vortrags. Wenn ich einen Mathematiker sehr hoch schätze und in dieser Hinsicht etwas vermisse, denke ich immer zuerst, dass es \textls{nur} an der Darstellung liegt, und nehme nur ungern an, dass der Gedanke selbst leer ist. Bei den drei Aufsätzen von Herrn \textsc{Ivory} bin ich leider dazu gezwungen.

Die Abhandlung über die kleinsten Quadrate ist wirklich unter aller Kritik. Welche Verworrenheit, Unklarheit und welcher völlige Mangel an logischer Schlüssigkeit! Darüber könnte man noch hinwegsehen, wenn tatsächlich ein wirklicher Gedanke zugrunde läge. Das ist aber keineswegs der Fall. Wie ist es möglich, dass man uns das Geschwätz auf S. 163 und 164 als einen \textls{proof} that this principle leads \textls{necessarily} to the method of the least squares anpreist!

Tatsächlich folgt daraus keineswegs die \textls{Notwendigkeit} der Methode der kleinsten Quadrate. Man kann das Ganze\HistoricalNote{*)}{*)}{Natürlich mit Ausnahme der Folgerung, dass $\varepsilon\varepsilon+\varepsilon'\varepsilon'+\varepsilon''\varepsilon''+\cdots$ ein Minimum wird. Aber genau das soll ja bewiesen werden und kein \textls{Axiom} sein.} auf \textls{jede beliebige} andere Behandlung der gegebenen Gleichungen anwenden. Man könnte sie zum Beispiel, statt sie mit $a$, $a'$, $a''$ zu multiplizieren, mit irgendeiner Potenz dieser Koeffizienten multiplizieren oder auch durch sie dividieren. Die Abhandlung scheint tatsächlich im Halbschlaf geschrieben worden zu sein. Sie ersparen es mir wohl, noch weitere einzelne Stellen herauszugreifen, die Ihnen sofort selbst auffallen werden, sofern Sie nicht von der vorgefassten Meinung ausgehen, dass irgendetwas daran ist.

Der andere Aufsatz über die Pendellängen scheint mir kaum besser. Die Einwände, die er gegen die Anwendung der Methode der kleinsten Quadrate erhebt, zeigen nur, dass ihr Geist Herrn I. völlig fremd ist. Ich finde, dass Herr \textsc{Sabine} diese Methode im Wesentlichen ganz richtig angewandt hat (— die Art der Anwendung ist allerdings nicht die zweckmäßigste, und unter uns gesagt hatte ich eigentlich auch \textsc{Sabine}, \textsc{Paucker} und \textsc{Muncke} in der neuen Ausgabe von \textsc{Gehler} bei meiner
